\begin{papercard}{shipp2020}
\cf{Stages} \stagetag{D} \stagetag{C} (also \stagetag{A}). A mixed-method audit of 30 Android menstruation apps. Policy content and language are compared with the data items each app actually requires, the third-party libraries embedded and the observed network traffic.

\cf{Research questions} Q1: how well do developers inform users in policies and privacy communications? Q2: is the information clear and understandable? Q3: do apps behave differently from what the policy states? \pc{shipp2020}{2}{how well do developers inform menstruapp users about their privacy}. Each method step is mapped to a question \pc{shipp2020}{3}{to see how well developers inform their}.

\cf{Sampling} Google Play, new Google account and new Android device, downloads in June 2018 \pc{shipp2020}{3}{using a new Google account on a new}. Five search terms: \emph{period tracker, period, menstrual cycle, menstruation tracker, menstrual calendar}. Results were ranked by downloads and ratings, apps under 100{,}000 downloads were filtered out, and the search continued until 30 apps (the study's workload limit) were reached \pc{shipp2020}{3}{The search terms used to find the apps were}. The rule follows Adhikari et al.\ (highest downloads and ratings) \pc{shipp2020}{3}{This method was in line with Adhikari et al}. Free apps only, because they reach wider audiences. Exceptions: two paid versions with $>$100k downloads were also analysed, because paid versions often inherit free-version practices, and Natural Cycles was included because its trial gives full functionality \pc{shipp2020}{3}{Our study was initially restricted to free apps}. The authors argue that 30 apps cover $>$90\,\% of downloads among $>$200 period apps \pc{shipp2020}{15}{our 30 app analysis is able to cover}.

\cf{Developer background} Desk research on each developer's business model, other products and marketing (Play Store, developer site, search engines), to see whether menstrual data is treated as sensitive or as ``just another data point'' \pc{shipp2020}{3}{We conducted background research on each of the}.

\cf{Policy location and in-app notices} Recorded whether the policy is available in the app, the store, or both. Each app was used manually for 10 minutes (account creation, period entry) while looking for any privacy notice or agreement request \pc{shipp2020}{4}{We manually interacted with each app for 10 minutes}.

\cf{Coding categories} \emph{Scope}, adopted from Sunyaev et al.: single app; multiple apps; back-end services; developer homepage; company-based services; no relation to app \pc{shipp2020}{4}{In our work we use the same categories}. \emph{Purpose}: app operation, personalisation, and secondary use split into scientific use, aggregation/analytics and third parties \pc{shipp2020}{4}{we found that the secondary use category can be split}. \emph{User rights}: access, deletion, portability as defined by the GDPR (chosen as the most tangible and common GDPR references), plus developer contact details and a named DPO \pc{shipp2020}{4}{described the user rights to data access, deletion and}. \emph{Data types}: fixed classes extending Sunyaev's scheme: PII, period data, reasons for use, sex logging, collection methods, test results, other health information, each with sub-items (Appendix B) \pc{shipp2020}{5}{These data classes fell into the following categories}. For each app and item the appendix marks Required / Optional / Not collected and whether the policy explicitly mentions it \pc{shipp2020}{19}{Types of information collected by the app and mentioned by the privacy policy}.

\cf{Language analysis} Operationalises obfuscation through three word classes: qualitative adjectives/adverbs that modify emphasis, temporal adverbs that blur frequency, and conditional verbs that introduce uncertainty \pc{shipp2020}{5}{Qualitative adjectives and adverbs that modify}. Procedure: spaCy tokenisation, POS extraction, manual filtering, extraction of sentences containing the retained words, and reading each word in context \pc{shipp2020}{5}{We used spaCy to split each privacy policy into tokens}. Verbs were counted only when the subject was the user, the developer, the app/company, ``we'' or ``you''. The full word lists are in Appendix C (e.g.\ verbs: could, would, may, might, can) \pc{shipp2020}{20}{For verbs, only those sentences having the user}. Flesch Reading Ease was computed per policy \pc{shipp2020}{5}{For each privacy policy we also calculated their}.

\cf{Behavioural comparison} Third-party libraries were detected with LibScout (library profiles from March 2019, partial matches discarded). For each library: was it disclosed, was the data described, and was the third-party policy linked? \pc{shipp2020}{6}{For this, we first used LibScout}. Network traffic: Android 6.0 with a user-installed certificate and a decryption proxy, a fresh install, about 10 minutes of use, distinctive data patterns entered to match in captures, and sync/backup triggered \pc{shipp2020}{6}{We installed a self-signed certificate and redirected all}. Certificate pinning blocked capture for three apps \pc{shipp2020}{6}{due to the use of certificate pinning}.

\cf{Key findings} 70\,\% linked the policy in both store and app. Four apps had no policy and no other privacy communication, although they had at least 11.5 million downloads. One policy link returned ``File Not Found'' \pc{shipp2020}{7}{There were 4 apps with no privacy policies}. Only 8 policies addressed the app specifically; 13 covered multiple apps or services \pc{shipp2020}{7}{In only 8 apps the respective privacy policy was directly}. Median Flesch 45, mean 47, i.e.\ university-level reading \pc{shipp2020}{10}{The median and average Flesch reading ease score ob}. ``May'' with the developer as subject was the main source of uncertainty \pc{shipp2020}{10}{obfuscation occurred due to the}. Required data exceeded the data mentioned in 56\,\% of apps. 28 apps required period data, 14 sent it to servers, and only 6 policies mentioned it \pc{shipp2020}{11}{The data required was greater than}. Only 5 policies named embedded libraries together with the data they receive \pc{shipp2020}{12}{explicitly mentioned the third party libraries that were}. Developers were contacted with the results \pc{shipp2020}{13}{We contacted the developers in October 2019 to share our}.

\cf{Limitations stated} Manual qualitative steps limit scalability. LibScout's library database could miss newer libraries, and some traffic was non-standard encoded. Apps and policies change over time, so policy hashes were archived on GitHub \pc{shipp2020}{15}{uploaded the privacy policies hashes of apps used to}.

\cf{Use for this thesis} Reusable: the scope categories (also in Sunyaev and Malki), the data-item matrix of required/optional/mentioned, the hashing of captured policy versions, and an operational, reproducible definition of hedging language (word lists, subject restriction, sentence-in-context reading). This is a direct precedent for coding provider \emph{stance} or hedged harm language. Not transferable: traffic and library analysis, which lie outside the communication-only scope. Flesch scores are English-specific and do not apply to German text without an adapted formula.
\end{papercard}
